\section*{Chapitre \ref{ch:s1:cross-asset-lead-lag} --- Avance–retard entre classes d'actifs}

\begin{solution}{exo:s1:cross-asset-lead-lag:1}
$\Phi^{-1}(1 - 0.05/2\,000) = 4.06$ pour 1\,000 tests~; $\Phi^{-1}(1 - 0.05/200) = 3.48$ pour 100.
\end{solution}

\begin{solution}{exo:s1:cross-asset-lead-lag:2}
Pour un retard, la covariance vaut $0.08\,\sigma_x^2/5$, donc la corrélation vaut $0.08 \times (25/8)/5 = 0.05$ et $t = 0.05\sqrt{1\,260} = 1.77$. Avec la
fenêtre de cinq jours, la covariance vaut $0.08\,\sigma_x^2$ et la volatilité de la fenêtre $\sigma_x\sqrt5$, donc la corrélation vaut $0.08 \times (25/8)/\sqrt5 = 0.112$
et $t = 0.112\sqrt{1\,255} = 3.96$, proche des 3.2 à 4.1 mesurés.
\end{solution}

\begin{solution}{exo:s1:cross-asset-lead-lag:3}
Il somme les produits de rendements sur toutes les paires d'intervalles qui se chevauchent, sans échantillonner les deux séries sur une grille commune~; il ne rejette aucun tick, et il ne
crée pas les rendements nuls qu'engendre l'échantillonnage au dernier tick, lequel biaise les corrélations vers zéro (effet Epps du livre~4).
\end{solution}

\begin{solution}{exo:s1:cross-asset-lead-lag:4}
$-0.51 + 1.0 = 0.49$ points de base avant coûts. Le prix observé du suiveur est son dernier tick, qui arrive à des instants aléatoires~: deux secondes après la
variation du meneur, certains ticks du suiveur n'ont pas encore reflété ce mouvement, de sorte qu'une petite part de celui-ci n'est pas encore intégrée dans les prix observés.
\end{solution}

\begin{solution}{exo:s1:cross-asset-lead-lag:5}
Le balayage par fenêtre regroupe les indices de cinq retards dans une seule statistique et ne coûte que 100 tests, si bien qu'il gagne en puissance pour une avance étalée. Le balayage retard par retard
teste chaque retard isolément contre un seuil plus strict, mais un retard qui se trouve fort dans l'échantillon (lien~1 à cinq jours, $t = 4.7$) peut le franchir
alors que sa statistique de fenêtre, diluée par des retards plus faibles, ne le franchit pas.
\end{solution}

\begin{solution}{exo:s1:cross-asset-lead-lag:6}
Il peut signaler les deux liens~: la matière première corrélée hérite d'une partie de la corrélation prédictive du vrai meneur. Régresser le suiveur sur les fenêtres des deux meneurs
ensemble (le coefficient du faux lien doit s'annuler), et privilégier les liens dotés d'un mécanisme.
\end{solution}

\begin{solution}{exo:s1:cross-asset-lead-lag:7}
Le balayage retard par retard trouve deux liens (lien~3 à un jour et lien~1 à cinq jours), le balayage par fenêtre les trois liens implantés ($t$ de 5.3, 6.1 et 6.2) et un
quatrième, matière première~10 vers industrie~4, qui est faux ($t = 3.7$)~: il repose en partie sur sa corrélation de 0.21 avec le vrai meneur de l'industrie~4 et
en partie sur du bruit. Le ratio de Sharpe hors échantillon (out-of-sample) atteint 2.56 après coûts~; le faux lien coûte peu.
\end{solution}

\begin{solution}{exo:s1:cross-asset-lead-lag:8}
30\,000 tests paire--retard, dont les 50 plus forts ont été retenus~: la plupart relèvent du bruit, et le backtest les a sélectionnés dans l'échantillon (in-sample). Corriger le nombre de tests, exiger
un mécanisme, estimer sur une période et négocier sur une autre, et imputer les coûts aux 50 paires négociées quotidiennement.
\end{solution}

\begin{solution}{pb:s1:cross-asset-lead-lag:1}
\begin{enumerate}
\item Une prévision d'un marché à partir du passé d'un autre~; l'information atteint les prix à des vitesses différentes parce que les investisseurs se spécialisent et suivent
peu de marchés.
\item Des secteurs comme le commerce de détail, les métaux et le pétrole prévoient le marché jusqu'à deux mois~; dans l'échantillon étendu, un noyau robuste le fait encore, mais
moins de secteurs, avec une variation dans le temps.
\item Les secteurs fournisseurs et clients liés se prédisent mutuellement, moins lorsque davantage d'investisseurs informés les suivent.
\item L'avance des grandes valeurs sur les petites se situe surtout à l'intérieur des secteurs et provient de la diffusion lente des mauvaises nouvelles.
\item 0.08 fois la moyenne des cinq derniers rendements quotidiens d'une matière première, dans trois liens~; une corrélation d'environ 0.05 par retard.
\item Retard par retard, 1\,000 tests à 4.06~; la fenêtre de cinq jours, 100 tests à 3.48.
\item Un lien à cinq jours ($t = 4.7$)~; les deux autres liens ($t = 4.1$ chacun), en laissant échapper le premier ($t = 3.2$)~; rien de faux.
\item Un ratio de Sharpe de 2.09 avant coûts et de 1.73 après.
\item Un suiveur qui voit le prix du meneur avec deux secondes de retard, les deux séries étant bruitées au niveau des ticks~; corrélation croisée de Hayashi--Yoshida.
\item À deux secondes, dans les vingt séances.
\item Négocier le suiveur après une variation du meneur de plus de 2 points de base en deux secondes, conserver deux secondes, un demi-point de base dans chaque sens.
\item 0.73, 0.35, $-0.51$ et $-0.81$ points de base pour des latences de 0 à 3 secondes.
\item \textbf{Résultat clé.} Quotidien~: une avance à cinq jours détectée sur deux des trois liens implantés par le balayage par fenêtre (un seul par le balayage retard par retard), négociée à un ratio de Sharpe
de 1.73 après coûts hors échantillon~; en infrajournalier, une avance de deux secondes détectée dans les vingt séances, valant 0.73 points de base par transaction sans latence, 0.35 à
une seconde, et une perte à partir de deux secondes.
\item Statistiques (peu de liens faibles parmi de nombreux candidats)~; vitesse et coûts.
\item Il fixe la forme du balayage et le nombre de tests, ce qui détermine la puissance.
\item Les traders les apprennent et les font disparaître en les exploitant, car la plupart des signaux publiés s'érodent.
\item Balayages glissants, échantillons scindés, et un lien désactivé en simulation pour voir à quelle vitesse le test le remarque.
\item L'avance intrajournalière du contrat à terme sur le comptant.
\item Le chapitre~17 a mesuré la part d'un mouvement d'actualité qui subsiste à chaque latence~; ici, l'avance d'un actif sur un autre constitue l'actualité.
\item Celui qui voit le meneur en premier et négocie le suiveur au coût le plus bas.
\end{enumerate}
\end{solution}

\begin{solution}{iq:s1:cross-asset-lead-lag:1}
Corréler les rendements du suiveur aux rendements passés du meneur aux horizons pertinents, dans les deux sens, avec un estimateur adapté aux données
(Hayashi--Yoshida pour les ticks)~; corriger le nombre de retards et de paires~; vérifier la stabilité dans le temps et hors échantillon~; exiger un mécanisme.
\end{solution}

\begin{solution}{iq:s1:cross-asset-lead-lag:2}
À 1~\% environ 100 sur 10\,000 seraient significatifs par hasard~; 40 est moins que le hasard n'en produit. Rien n'a été trouvé~; appliquer une correction pour tests multiples et exiger
un mécanisme avant de chercher à nouveau.
\end{solution}

\begin{solution}{iq:s1:cross-asset-lead-lag:3}
Synchroniser les horloges (utiliser les horodatages de la bourse, ou mesurer et corriger le retard de chaque flux), garder les ticks asynchrones et appliquer l'estimateur de
Hayashi--Yoshida en corrélation croisée sur une grille de retards, et lire l'avance au sommet ou au centre de symétrie de la fonction de corrélation.
\end{solution}

\begin{solution}{iq:s1:cross-asset-lead-lag:4}
Les concurrents sont devenus plus rapides, les fournisseurs de liquidité du suiveur ont commencé à coter en fonction du meneur, les frais ou les fourchettes ont changé, ou la latence propre de la société a augmenté~;
mesurer de nouveau l'avance et la latence de la société.
\end{solution}

\begin{solution}{iq:s1:cross-asset-lead-lag:5}
Le lien peut se rompre (changement de régime, concurrents plus rapides), le suiveur peut connaître un saut de cours sur sa propre actualité, et la position reste non couverte contre le
risque spécifique du suiveur~; en infrajournalier, des données périmées ou un flux lent transforment la transaction en sélection adverse.
\end{solution}

\begin{solution}{iq:s1:cross-asset-lead-lag:6}
$\operatorname{Cov}(y_t, x_{t-k}) = \beta\sigma_x^2/L$, donc $\rho_1 = \beta\sigma_x/(L\sigma_y)$. La fenêtre $W = \sum_k x_{t-k}$ a une variance $L\sigma_x^2$ et
$\operatorname{Cov}(y_t, W) = \beta\sigma_x^2$, donc $\rho_W = \beta\sigma_x/(\sqrt{L}\sigma_y)$. Sur le même échantillon, le rapport des statistiques t vaut
$\rho_W/\rho_1 = \sqrt{L}$~: la fenêtre a $\sqrt{L}$ fois la puissance de tout retard isolé.
\end{solution}
